\BourbakiExerciseGroup

\begin{enumerate}
\def\labelenumi{\arabic{enumi})}
\tightlist
\item
  设 \(U = (\alpha_{ij})\) 是主理想整环 A 上一个 \(m\) 行 \(n\) 列的矩阵。a) 首先假设诸 \(\alpha_{ij}\) 整体互素。证明存在两个元素属于 A 的可逆方阵 \(P\) 和 \(Q\)，使得矩阵 \(PUQ\) 的某一元素等于 1。（对每个 \(\alpha_{ij} \neq 0\) ，设 \(S(\alpha_{ij})\) 为 \(\alpha_{ij}\) 分解为不可约元时诸指数之和，并设 \(s(U)\) 为诸数 \(s(\alpha_{ij})\) 中最小者；若 \(s(U) > 0\) ，证明存在两个可逆方阵 \(R\) 和 \(S\)，使得
\end{enumerate}

\[
s (R U S) <   s (U);
\]

利用贝祖恒等式，并在必要时对行和列作置换后，把所考虑的矩阵 \(R\) 和 \(S\) 限于形如 \(\begin{pmatrix} T & 0 \\ 0 & I \end{pmatrix}\) 的矩阵，其中 \(T\) 是一个 \(2 \times 2\) 矩阵，\(I\) 是单位矩阵，或限于此类矩阵与置换矩阵的乘积。）

\begin{enumerate}
\def\labelenumi{\alph{enumi})}
\setcounter{enumi}{1}
\tightlist
\item
  若 \(\delta_1\) 是 \(U\) 的诸元素的一个最大公因子，证明存在两个可逆矩阵 \(P_1, Q_1\)，使得
\end{enumerate}

\[
P _ {1} U Q _ {1} = \left( \begin{array}{c c} \delta_ {1} & 0 \\ 0 & U _ {1} \end{array} \right)
\]

其中 \(U_{1}\) 的所有元素都可被 \(\delta_{1}\) 整除（利用 a））。

\begin{enumerate}
\def\labelenumi{\alph{enumi})}
\setcounter{enumi}{2}
\tightlist
\item
  对于 \(\mathbf{A} = \mathbf{Z}\) ，利用 \(b\) 得到一个计算元素属于 \(\mathbf{Z}\) 的具体矩阵的不变因子的方法。将此方法应用于矩阵
\end{enumerate}

\[
\left( \begin{array}{c c c c} 6 & 8 & 4 & 20 \\ 12 & 12 & 18 & 30 \\ 18 & 4 & 4 & 10 \end{array} \right).
\]

\begin{enumerate}
\def\labelenumi{\arabic{enumi})}
\setcounter{enumi}{1}
\item
  设 A 为主理想整环。则 \(A^{q}\) 的两个子模 M 与 N 在 \(A^{q}\) 的自同构下互为变换，当且仅当它们关于 \(A^{q}\) 具有相同的不变因子。
\item
  设 A 为主理想整环，其分式域为 K，设 E 为 K 上的向量空间，设 M 为包含于 E 中的秩为 n 的有限生成 A-模，并设 N 为包含于 KM 中的秩为 \(p \leqslant n\) 的有限生成 A-模。证明存在 M 的一组基 \((e_{i})_{1 \leqslant i \leqslant n}\) 以及 N 的一组由向量 \(\alpha_{i}e_{i}\) ( \(1 \leqslant i \leqslant p\) ) 组成的基，其中 \(\alpha_{i} \in K\) ，且 \(\alpha_{i}\) （关于环 A）整除 \(\alpha_{i+1}\) ，对 \(1 \leqslant i \leqslant p - 1\) ；同时证明分式理想 \(A\alpha_{i}\) 是唯一确定的（即 N 关于 M 的不变因子）。
\item
  设 \(G\) 是一个有限阿贝尔群。
\end{enumerate}

\begin{enumerate}
\def\labelenumi{\alph{enumi})}
\item
  对每个素数 \(p\) ，证明阶为 \(p\) 的元素在 \(G\) 中的个数是 \(p^N - 1\) ，其中 \(N = \sum_{n=1}^{\infty} m(p^n)\) （采用 VII，第 24 页，命题 9 的记号）。
\item
  由此推出，若方程 \(x^p = 1\) 对每个素数 \(p\) 在 \(G\) 中至多有 \(p\) 个解，则 \(G\) 是循环群。
\end{enumerate}

\begin{enumerate}
\def\labelenumi{\arabic{enumi})}
\setcounter{enumi}{4}
\item
  设 \(G\) 为一个 \(n\) 阶有限交换群。证明：只要整数 \(q\) 整除 \(n\) ，就存在一个 \(q\) 阶子群包含在 \(G\) 中（把 \(G\) 分解为不可分解循环子群的直和）。
\item
  设 A 是一个交换环，并设 E 是一个 A-模。设 \(\mathcal{E}\) 表示（无算子的）阿贝尔群 E 的自同态环 \(\operatorname{End}_{\mathbf{z}}(\mathbf{E})\) ；对每个子环 \(\mathbf{B} \subset \mathcal{E}\) ，设 \(\mathbf{B}'\) 表示 B 在 \(\mathcal{E}\) 中的中心化子（即 \(\mathcal{E}\) 中由与 B 的每个元素都可交换的元素组成的子环）。
\end{enumerate}

\begin{enumerate}
\def\labelenumi{\alph{enumi})}
\item
  设 \(A_{0}\) 为 E 的由位似变换 \(x \mapsto \lambda x (\lambda \in A)\) 组成的子环；若 a 是 E 的零化理想，则已知 \(A_{0}\) 同构于 A/a，且它在 E 中的中心化子 \(A_{0}^{\prime}\) 是 E 作为 A-模的自同态环 \(\operatorname{End}_{\mathbf{A}}(\mathbf{E})\) 。证明中心化子 \(A_{0}^{\prime\prime}\) （\(A_{0}^{\prime}\) 的中心化子）是交换的，且其中心化子 \(A_{0}^{\prime\prime\prime}\) 等于 \(A_{0}^{\prime}\) （注意 \(A_{0} \subset A_{0}^{\prime}\) ）。
\item
  如果子模 F 是 E 的一个直因子，证明 \(u(F) \subset F\) 对所有 \(u \in A_{0}''\) （考虑由 E 分解为 F 与另一个子模的直和所得到的从 E 到 F 的投影）。
\item
  假设 E 是循环子模的一个族 \((F_{i})\) 的直和。若 \(u \in A_{0}''\) ，证明对每个指标 i，存在元素 \(\alpha_{i} \in A\) ，使得 \(u(x) = \alpha_{i}x\) 对所有 \(x \in F_{i}\) 成立（利用 b））。
\item
  假设 E 是循环子模的一个序列 \((\mathrm{E}_{n})\) （有限或无限）的直和，使得：若 \(b_{n}\) 表示 \(E_{n}\) 的零化理想，则序列 \((\mathfrak{b}_{n})\) 是递增的。证明 \(A_{0}^{\prime\prime}=A_{0}\) （利用 c)，以及对每个指标 i\textgreater1 存在一个自同态 \(v_{i}\in A_{0}^{\prime}\) 把 \(E_{i-1}\) 映入 \(E_{i}\) 这一事实。）
\end{enumerate}

将此应用于 A 为主理想整环且 E 为有限生成 A-模的情形。

\begin{enumerate}
\def\labelenumi{\arabic{enumi})}
\setcounter{enumi}{6}
\tightlist
\item
  设 A 为主理想整环，其分式域为 K，并设 E 为秩 n 的有限生成无挠 A-模。设 \(E^{*}\) 为 E 的对偶；则 E 与 \(E^{*}\) 均同构于 \(A^{n}\) .
\end{enumerate}

\begin{enumerate}
\def\labelenumi{\alph{enumi})}
\item
  设 M 为 E 的子模；证明：若 \(M^{0}\) 是 \(E^{*}\) 中与 M 正交的子模，则模 \(E^{*}/M^{0}\) 是无挠的，从而 \(M^{0}\) 是 \(E^{*}\) 的直因子；子模 \(M^{00}\) （E 中与 \(M^{0}\) 正交者）等于 \(E \cap KM\) （把 E 视为典范地嵌入到 K 上一个 n 维向量空间中）。可把 \(E^{*}/M^{0}\) 视为典范地嵌入到 M 的对偶 \(M^{*}\) 中；证明 \(E^{*}/M^{0}\) 关于 \(M^{*}\) 的不变因子等于 M 关于 E 的不变因子（利用 VII 第 18 页定理 1）。
\item
  设 F 是第二个有限生成的无挠 A-模，并设 u 是从 E 到 F 的线性变换。证明 \(u(F^{*})\) 关于 \(E^{*}\) 的不变因子与 \(u(E)\) 关于 F 的那些不变因子相等（按 VII 第 21 页命题 5 的方式进行）。
\end{enumerate}

\begin{enumerate}
\def\labelenumi{\arabic{enumi})}
\setcounter{enumi}{7}
\tightlist
\item
  设 G 是主理想整环 A 上的一个有限长度模。对于满足 \(N \subset M\) 且使 G 同构于 M/N 的任意一对自由模 M 和 N，N 关于 M 的不变因子中异于 A 的那些与 G 的不变因子相同（VII，第 20 页，定义 2）；它们的个数称为 G 的秩。
\end{enumerate}

\begin{enumerate}
\def\labelenumi{\alph{enumi})}
\item
  证明 G 的秩 r 是使 G 为其直和的循环子模的最少个数，且等于 G 的诸 \(\pi\) -准素分量的秩中的最大者。b) 证明 G 的子模或商模的秩至多等于 G 的秩（把 G 视为两个秩为 r 的自由模的商）。
\item
  对每个 \(\lambda \in \mathbf{A}\) ，证明子模 \(\lambda G\) 的秩等于 \(G\) 的不变因子中不整除 \(\lambda\) 的那些的个数（注意，若 \(E = A / (A\alpha)\) ，则模 \(\lambda E\) 同构于 \((A\lambda) / ((A\alpha) \cap (A\lambda))\) ）。由此推出，第 \(k\) 个不变因子（按递减顺序，属于 \(G\) ）是那些 \(\lambda \in A\) 的一个最大公因子，其中 \(\lambda G\) 的秩 \(\leqslant r - k\) 。
\item
  设 \(A\alpha_{k}\) ( \(1 \leqslant k \leqslant r\) ) 是 G 的不变因子，按递减顺序排列。设 H 是 G 的秩为 r - q 的子模，并设 \(A\beta_{k}\) ( \(1 \leqslant k \leqslant r - q\) ) 是它的不变因子，按递减顺序排列。由 b) 和 c) 推出 \(\beta_{k}\) 整除 \(\alpha_{k+q}\) ，对 \(1 \leqslant k \leqslant r - q\) 。类似地证明，若 G/H 的秩为 r - p 且不变因子为 \(A\gamma_{k}\) ( \(1 \leqslant k \leqslant r - p\) )（按递减顺序），则 \(\gamma_{k}\) 整除 \(\alpha_{k+p}\) ，对 \(1 \leqslant k \leqslant r - p\) 。
\item
  反之，设 L 是秩为 r - q 的挠模，其不变因子 \(A\lambda_{k}\) ( \(1 \leqslant k \leqslant r - q\) ) 满足 \(\lambda_{k}\) 整除 \(\alpha_{k+q}\) ，对 \(1 \leqslant k \leqslant r - q\) 。证明存在 G 的子模 M 和 N，使 L 同构于 M 且同构于 G/N（把 G 分解为同构于 A/A \(\alpha_{k}\) 的子模的直和）。
\end{enumerate}

\begin{enumerate}
\def\labelenumi{\arabic{enumi})}
\setcounter{enumi}{8}
\tightlist
\item
  设 A 是一个主理想整环，K 为其分式域；设 E 是 K 上的向量空间，设 M 是包含在 E 中的秩为 n 的有限生成 A-模，设 N 是包含在 KM 中的秩为 p 的有限生成 A-模，设 P 是包含在 KN 中的秩为 q 的有限生成 A-模。
\end{enumerate}

\begin{enumerate}
\def\labelenumi{\alph{enumi})}
\item
  设 \(\mathbf{A}\alpha_{i}\) ( \(1 \leqslant i \leqslant p\) ) 是 \(\mathbf{N}\) 关于 \(\mathbf{M}\) 的不变因子，按递减顺序排列（习题 3）。证明，对 \(1 \leqslant k \leqslant p\) ，元素 \(\alpha_{k}\) 是那些元素 \(\lambda \in \mathbf{K}\) 的一个最大公因子，这些元素使商模 \((\mathbf{N} + \lambda (\mathbf{KN} \cap \mathbf{M})) / \mathbf{N}\) 的秩 \(\leqslant p - k\)（参见习题 8，d））。
\item
  设 \(\mathbf{A}\beta_{j}\) ( \(1 \leqslant j \leqslant q\) ) 是 \(\mathbf{P}\) 关于 \(\mathbf{M}\) 的不变因子，\(\mathbf{A}\gamma_{j}\) ( \(1 \leqslant j \leqslant q\) ) 是 \(\mathbf{P}\) 关于 \(\mathbf{N}\) 的不变因子，均按递减顺序排列。证明 \(\alpha_{1}\gamma_{j}\) 整除 \(\beta_{j}\) ，对 \(1 \leqslant j \leqslant q\) （参见习题 8）。
\item
  证明 \(\gamma_{1}\alpha_{j}\) 整除 \(\beta_{j}\) ，对 \(1 \leqslant j \leqslant q\) 。（首先考虑 \(n = p = q\) 的情形，并应用习题 8。一般地，证明我们总可以假设 \(p = n\) ，并考虑 M 的一个子模 Q，使得 \(\mathbf{P} + \rho \mathbf{Q}\) 的秩为 \(p\) ，对所有 \(\rho \in \mathbf{A}\) ；然后选取足够小的理想 \(\mathbf{A}\rho\) ，并应用 VII 第 20 页命题 4。）
\end{enumerate}

\(d\) ) 对 N 的任意秩为 \(k \leqslant p\) 的子模 H，设 \((\mu_{\mathrm{H}})\) 为由那些 \(\mu \in \mathbf{K}\) 组成的分式理想，它们满足 \(\mu (\mathbf{M} \cap \mathbf{KH}) \subset \mathbf{H}\) 。由 \(c\) ) 推出，对每个 \(1 \leqslant k \leqslant p\) ，元素 \(\alpha_{k}\) 是诸 \(\mu_{\mathrm{H}}\) 的一个最大公因子，这里 H 遍历 N 的所有秩为 \(k\) 的子模。

\begin{enumerate}
\def\labelenumi{\arabic{enumi})}
\setcounter{enumi}{9}
\tightlist
\item
  设 A 为主理想整环，M 是 \(E = A^{n}\) 中秩为 n 的子模；设 \(A\alpha_{i} (1 \leqslant i \leqslant n)\) 为 M 关于 E 的不变因子，按递减顺序排列。设 N 是 M 的在 M 中有补的子模（即满足 \(N = M \cap KN\) ，其中 K 是 A 的分式域）。
\end{enumerate}

\begin{enumerate}
\def\labelenumi{\alph{enumi})}
\item
  设 P 是 N 在 M 中的一个补。证明 E/M 的子模 \((P + (E \cap KN))/M\) 同构于 \((E \cap KN)/N\) 。由此推出，若 N 的秩为 p 且 \(A\beta_{k}\) ( \(1 \leqslant k \leqslant p\) ) 是 N 关于 E 的不变因子，则 \(\beta_{k}\) 整除 \(\alpha_{k+n-p}\) ，对 \(1 \leqslant k \leqslant p\) （利用习题 8，d））。
\item
  给出一个例子，说明模 E/M 不必同构于模 \((E \cap KN)/N\) 与 \(E/(P + (E \cap KN))\) 的乘积（取 E/M 为不可分解的）。
\end{enumerate}

\begin{enumerate}
\def\labelenumi{\arabic{enumi})}
\setcounter{enumi}{10}
\item
  设 A 为主理想整环，H 是 \(E = A^{n}\) 的一个子模，\(H^{0}\) 是 E 的对偶 \(E^{*}\) 中与 H 正交的子模（习题 7）。设 M 是 E 的一个秩为 p 的子模，并设 \(A\alpha_{i}\) ( \(1 \leqslant i \leqslant p\) ) 为 M 关于 E 的不变因子，按递减顺序排列。设 \((\nu_{\mathrm{H}})\) 为 A 的诸理想 \(f(\mathbf{M})\) 的最大公因子，这里 f 遍历 \(H^{0}\) ；证明，对 \(1 \leqslant k \leqslant p\) ，元素 \(\alpha_{k}\) 是 E 的秩为 k - 1 且在 E 中有补的所有子模 H 所对应的 \(\nu_{H}\) 的一个最大公因子（首先考虑 k = 1 的情形；把结果应用于 E/H 的子模 \((\mathbf{M} + \mathbf{H})/\mathbf{H}\) ，并对 E/M 的商模 \(E/(M + H)\) 利用习题 8，d））。
\item
  设 G 是环 A 上的有限长度模。证明：G 的子模 H 若与 G 同构，则必有 H = G（参见 II，第 212 页，命题 16）。
\item
  设 A 为主理想整环，其分式域为 K。
\end{enumerate}

\begin{enumerate}
\def\labelenumi{\alph{enumi})}
\item
  证明：对于每个A-模M，典范同态 \(c_{\mathbf{M}}: \mathbf{M} \to \mathbf{D}(\mathbf{D}(\mathbf{M}))\) 是单射的（参见II，第389页，习题13）。
\item
  若 \(u: \mathbf{M} \to \mathbf{N}\) 是 A-模的一个同态，定义从 \(D(\operatorname{Im}(u))\) 到 \(\operatorname{Im}(D(u))\) 、从 \(D(\operatorname{Ker}(u))\) 到 \(\operatorname{Coker}(D(u))\) 以及从 \(D(\operatorname{Coker}(u))\) 到 \(\operatorname{Ker}(D(u))\) 的自然同构。
\item
  设 M 是一个 A-模，N 是 M 的一个子模。在第 9 节的记号下，证明 \(N^{00} = N\) （利用 a) 和 b)）。模 \(N^{0}\) 具有有限长度，当且仅当 M/N 具有有限长度。反之，对 D(M) 的任意有限长度子模 \(N'\) ，模 \(N'^{0}\) 是 M 的一个子模，使 \(M/N^{0}\) 具有有限长度，且 \(N' = N'^{00}\) 。
\end{enumerate}

\(d\) ) 若 \(\mathbf{M}\) 是挠 A-模，且 \(\mathbf{M}_{\pi}\) 表示 \(\pi\) -准素分量（属于 \(\mathbf{M}\) ），证明 \(\mathrm{D}(\mathbf{M})\) 同构于 \(\prod_{\pi \in \mathbb{P}}\mathrm{Hom}(\mathbf{M}_{\pi},\mathbf{U}_{\pi})\) ，其中 \(\mathbf{P}\) 是 A 的不可约元素的一个代表系，且 \(U_{\pi}\) 是 K/A 的 \(\pi\) -准素分量，对每个 \(\pi \in P\) （VII，第 54 页，习题 3）。

\begin{enumerate}
\def\labelenumi{\alph{enumi})}
\setcounter{enumi}{4}
\tightlist
\item
  由d)推出， \(\mathrm{D}(\mathrm{K}/\mathrm{A})\) 同构于乘积 \(\prod_{\pi\in P}\hat{\mathbf{A}}_{(\pi)}\) ，其中
\end{enumerate}

\(\hat{\mathbf{A}}_{(\pi)}\) 表示 VII，第 54 页，习题 2 中定义的环 \(\mathbf{A}_{(\pi)}\) 的完备化，其拓扑由取理想 \(\pi^n\mathbf{A}_{(\pi)}\) 为 0 的开邻域基来定义（此完备化还可等同于逆向极限 \(\lim (\mathbf{A} / \mathbf{A}\pi^n)\) ）。

\begin{enumerate}
\def\labelenumi{\alph{enumi})}
\setcounter{enumi}{5}
\item
  若 M 是有限长度的 A-模，且 a 是其零化子，证明 D(M) 同构于与 M 相伴的忠实 (A/a)-模的对偶。
\item
  设 M 是有限长度的 A-模。证明：若 N 是 M 的子模，则存在 M 的子模 P，使得 M/N 同构于 P，且 M/P 同构于 N（利用 VII，第 26 页，命题 10 及 VII，第 27 页，命题 12）；具有这些性质的 M 的两个子模称为互反的。证明：对所有 \(\alpha \in A\) ，M 的子模 \(M(\alpha)\) （由所有满足 \(x \in M\) 且 \(\alpha x = 0\) 的元素组成）与子模 \(\alpha M\) 是互反的。
\end{enumerate}

¶ 14) 设 M 是主理想整环 A 上的有限长度模，A 的分式域为 K。

\begin{enumerate}
\def\labelenumi{\alph{enumi})}
\tightlist
\item
  证明：若 N 是 M 的子模，则 M 的秩至多等于 N 与 M/N 的秩之和（令 M = E/H，其中 E 是秩为 n 的自由模，H 是 E 的秩为 n 的子模；若 N = L/H，其中 H ⊆ L ⊆ E，注意，若 N 的秩为 p，则存在 H 的秩为 n - p 的子模 R，使得 R 在 L 中有补 S，且 S ∩ H 是 R 在 H 中的补；此外，(E ∩ KR)/R 的秩 ≥ n - p，且同构于 E/L 的一个商模（习题 10，a））；利用习题 8，b）完成证明。 b) 设 N 是 M 的秩为 p 的子模，q 是 M/N 的秩。设 Aα \(_{i}\) （1 ≤ i ≤ n）为 M 的不变因子，Aβ \(_{k}\) （1 ≤ k ≤ p）为 N 的不变因子，均按递减顺序排列。证明 β \(_{k}\) 是 α \(_{k- (p + q - n)}\) （利用a）及习题8，c），注意到(λM)/(λN)同构于M/N的一个商模）。
\end{enumerate}

\begin{enumerate}
\def\labelenumi{\arabic{enumi})}
\setcounter{enumi}{14}
\tightlist
\item
  \begin{enumerate}
  \def\labelenumii{\alph{enumii})}
  \tightlist
  \item
    设 M 是主理想整环 A 上秩为 n 的有限长度模，并设 \(\mathbf{A}\alpha_{i}\) ( \(1 \leqslant i \leqslant n\) ) 为其按递减顺序排列的不变因子。设 \((\beta_{i})\) ( \(1 \leqslant i \leqslant n\) ) 是 A 中元素的一个序列，使得 \(\beta_{i}\) 整除 \(\alpha_{i}\) ，对 \(1 \leqslant i \leqslant n\) ，且 \(\alpha_{i}\) 整除 \(\beta_{i+1}\) ，对 \(1 \leqslant i \leqslant n-1\) 。设 \((\mathbf{M}_{i})_{1 \leqslant i \leqslant n}\) 是 M 的子模的一个序列，使得 M 是诸 \(M_{i}\) 的直和，且每个 \(M_{i}\) 同构于 \(A/A\alpha_{i}\) ；设 \(a_{i}\) 是 \(M_{i}\) 的一个生成元，对 \(1 \leqslant i \leqslant n\) 。令
  \end{enumerate}
\end{enumerate}

\[
b = \beta_ {1} a _ {1} + \frac {\beta_ {1} \beta_ {2}}{\alpha_ {1}} a _ {2} + \frac {\beta_ {1} \beta_ {2} \beta_ {3}}{\alpha_ {1} \alpha_ {2}} a _ {3} + \dots + \frac {\beta_ {1} \beta_ {2} \dots \beta_ {n}}{\alpha_ {1} \alpha_ {2} \dots \alpha_ {n - 1}} a _ {n}.
\]

证明 M 关于循环子模 Ab 的商模同构于 n 个模 \(A/A\beta_{i}\) ( \(1 \leq i \leq n\) ) 的直和。

\begin{enumerate}
\def\labelenumi{\alph{enumi})}
\setcounter{enumi}{1}
\tightlist
\item
  设 M 和 N 是 A 上两个有限长度模；设 \(A\alpha_{i}\) ( \(1 \leqslant i \leqslant n\) ) 和 \(A\beta_{k}\) ( \(1 \leqslant k \leqslant p\) ) 分别是 M 和 N 的不变因子，按递减顺序排列；假设 \(p \leqslant n\) ，\(\beta_{k}\) 整除 \(\alpha_{k+n-p}\) ，对 \(1 \leqslant k \leqslant p\) ，且 \(\alpha_{k}\) 整除 \(\beta_{k+(p+q-n)}\) ，对 \(1 \leqslant k \leqslant n-q\) ，其中 q 是一个整数，使得
\end{enumerate}

\[
q \leqslant n \leqslant p + q.
\]

证明存在 M 的一个与 N 同构的子模 P，使得 M/P 的秩 \(\leq q\) （把两个模 M 和 N 各自分解为 q 个模的直和，从而把问题归结为 q = 1 的情形；然后利用 a) 以及习题 13 的 g)）。

\begin{enumerate}
\def\labelenumi{\alph{enumi})}
\setcounter{enumi}{2}
\tightlist
\item
  对于 A 上的模 N，要使它与有限长度模 M 的一个直和因子同构，当且仅当 N 的每个初等因子都是 M 的初等因子，且其在 N 中的重数至多等于其在 M 中的重数。由此给出一个有限长度模 M 及其子模 N 的例子，使得 M 的秩等于 N 与 M/N 的秩之和，但 N 在 M 中没有补（利用 a)）。
\end{enumerate}

\begin{enumerate}
\def\labelenumi{\arabic{enumi})}
\setcounter{enumi}{15}
\tightlist
\item
  设 M 是交换环 A 上的一个模。若对 M 的每个自同态 u 都有 \(u(N) \subset N\) ，则 M 的子模 N 称为特征子模。对所有 \(\alpha \in A\) ，子模 \(M(\alpha)\) （由被 \(\alpha\) 零化的元素组成）与子模 \(\alpha M\) 都是特征的。
\end{enumerate}

\begin{enumerate}
\def\labelenumi{\alph{enumi})}
\item
  若 N 是 M 的特征子模，且 P 与 Q 是 M 的两个互补子模，证明 N 是 \(N \cap P\) 与 \(N \cap Q\) 的直和（考虑从 M 到 P 和 Q 上的投影）。
\item
  设 M 是主理想整环 A 上的一个有限长度模；设 \(A\alpha_{i}\) 为 M 的不变因子，按递减顺序排列 \((1 \leqslant i \leqslant n)\) ；设 \((\mathbf{M}_{i})_{1 \leqslant i \leqslant n}\) 是 M 的子模的一个序列，使得 M 是诸 \(M_{i}\) 的直和，且每个 \(M_{i}\) 同构于 \(A/A\alpha_{i}\) 。设 N 是 M 的一个特征子模，并设 \(A\beta_{i}\) 为 \(N \cap M_{i}\) 的零化理想 \((1 \leqslant i \leqslant n)\) ；若 \(\alpha_{i} = \beta_{i}\gamma_{i}\) ，证明 \(\beta_{i}\) 整除 \(\beta_{i+1}\) 且 \(\gamma_{i}\) 整除 \(\gamma_{i+1}\) ，对所有 \(1 \leqslant i \leqslant n-1\) （注意，对 \(1 \leqslant i \leqslant n-1\) ，存在 M 的一个把 \(M_{i+1}\) 映到 \(M_{i}\) 上的自同态，以及 M 的一个把 \(M_{i}\) 映到 \(M_{i+1}\) 的一个子模上的自同态）。
\item
  反之，设 N 是 M 的一个子模，它是子模 \(N_{i} \subset M_{i}\) 的直和；假设诸零化理想 \(A\beta_{i}\) （对应诸 \(N_{i}\) ， \(1 \leqslant i \leqslant n\) ）满足上述条件。证明 N 是 M 的一个特征子模。由此推出，若 M 的两个特征子模具有相同的不变因子，则它们相等。
\end{enumerate}

\begin{enumerate}
\def\labelenumi{\arabic{enumi})}
\setcounter{enumi}{16}
\tightlist
\item
  设 A 为主理想整环，\(\alpha\) 是 A 的一个非零元素，并设 E 为模 A/A \(\alpha\) 。则乘积 \(E^{n}\) 同构于 \(A^{n}/(\alpha A^{n})\) 。
\end{enumerate}

\begin{enumerate}
\def\labelenumi{\alph{enumi})}
\item
  设 \(u\) 是从 \(\mathbf{E}^n\) 到 \(\mathbf{E}^m\) 的一个线性映射；证明 \(u\) 可由一个线性映射 \(v\) （从 \(\mathbf{A}^n\) 到 \(\mathbf{A}^m\) ）通过过渡到商而得到。
\item
  设 \(A\beta_{i}\) ( \(1 \leqslant i \leqslant p \leqslant \inf(m, n)\) ) 是 \(v(A^{n})\) 关于 \(A^{m}\) 的不变因子，按递减顺序排列；设 \(q \leqslant p\) 为使得 \(\alpha\) 不整除 \(\beta_{i}\) 的指标 i 中最大者，并设 \(\delta_{i}\) 为 \(\alpha\) 与 \(\beta_{i}\) 的一个最大公因子，对 \(1 \leqslant i \leqslant q\) 。证明 u 的核 \(u^{-1}(0)\) 是 n - q 个同构于 E 的模与 q 个分别同构于 \(A/A\delta_{i}\) （ \(1 \leqslant i \leqslant q\) ）的循环模的直和。若 \(\alpha = \gamma_{i}\delta_{i}\) ( \(1 \leqslant i \leqslant q\) )，证明模 \(u(E^{n})\) 具有不变因子 \(A\gamma_{i}\) 。
\end{enumerate}

\begin{enumerate}
\def\labelenumi{\arabic{enumi})}
\setcounter{enumi}{17}
\tightlist
\item
  设 \(\mathbf{C}\) 为一个主理想整环。考虑线性方程组
\end{enumerate}

\[
\sum_ {j = 1} ^ {n} \alpha_ {i j} \xi_ {j} = \beta_ {i} (1 \leqslant i \leqslant m)
\]

其中系数与右端项都属于 C，设 A 表示 m 行 n 列的矩阵 \((\alpha_{ij})\) ，而 B 表示通过添加第 \((n+1)\) 列 \((\beta_{i})\) 扩充 A 而得到的矩阵。证明该方程组在 \(C^{n}\) 中至少有一个解，当且仅当下列条件满足：1) A 与 B 具有相同的秩 p；2) A 的 p 阶子式的一个最大公因子等于 B 的 p 阶子式的一个最大公因子（利用 VII 第 20 页命题 4 以及习题 9，c））。

\begin{enumerate}
\def\labelenumi{\arabic{enumi})}
\setcounter{enumi}{18}
\tightlist
\item
  设 \(\mathbf{C}\) 是一个主理想整环。考虑由 \(m\) 个«线性同余式»组成的方程组 \(\sum_{j=1}^{n} \alpha_{ij} \xi_j \equiv \beta_i (\bmod \alpha)\) ，其中系数、右端项以及元素 \(\alpha\) 都属于 C，且 \(\alpha\) 非零；设 \(A\) 表示矩阵 \((\alpha_{ij})\) ，设 \(B\) 表示通过添加列 \((\beta_i)\) 扩充 A 而得到的矩阵，并设 \(\delta_k\) （相应地 \(\delta_k'\) ) 为 \(k\) 阶子式（\(A\) 的，相应地 \(B\) 的）的一个最大公因子。则该方程组在 \(\mathbf{C}^n\) 中至少有一个解，当且仅当：\(a)\) 若 \(m \leqslant n\) ，\(\alpha^m\) , \(\alpha^{m-1}\delta_1\) , \ldots, \(\alpha\delta_{m-1}\) , \(\delta_m\) 的一个最大公因子等于 \(\alpha^m\) , \(\alpha^{m-1}\delta_1'\) , \ldots, \(\alpha\delta_{m-1}'\) , \(\delta_m'\) 的一个最大公因子；
\end{enumerate}

\begin{enumerate}
\def\labelenumi{\alph{enumi})}
\setcounter{enumi}{1}
\tightlist
\item
  若 \(m > n\) ，\(\alpha^{n+1}\) , \(\alpha^n \delta_1, \ldots, \alpha \delta_n\) 的一个最大公因子等于 \(\alpha^{n+1}\) , \(\alpha^n \delta_1', \ldots, \delta_{n+1}'\) 的一个最大公因子。
\end{enumerate}

将线性同余式方程组化为线性方程组，并利用习题 18。

\begin{enumerate}
\def\labelenumi{\arabic{enumi})}
\setcounter{enumi}{19}
\tightlist
\item
  \begin{enumerate}
  \def\labelenumii{\alph{enumii})}
  \tightlist
  \item
    设 \(\mathbf{C}\) 是一个主理想整环，并且设
  \end{enumerate}
\end{enumerate}

\[
f (x _ {1},.., x _ {p}) = \sum_ {(i _ {k})} \alpha_ {i _ {1} \dots i _ {p}} \xi_ {1, i _ {1}} \dots \xi_ {p, i _ {p}}
\]

是一个 \(p\) -线性形式，定义在 \(\mathbf{E}^p\) 上，其中 \(\mathbf{E}\) 是 C-模 \(\mathbf{C}^n\) （而 \(\xi_{ij}\) 表示第 \(j\) 个坐标，对应于 \(x_i\) ）。证明若 \(\delta\) 是诸系数 \(\alpha_{i_1\dots i_p}\) 的一个最大公因子，则存在一个元素 \((a_k)\in \mathbf{E}^p\) ，使得 \(f(a_1,\dots,a_p) = \delta\) （化归到 \(\delta = 1\) 的情形，并对 \(p\) 作归纳，利用 VII，第 18 页，定理 1）。

\begin{enumerate}
\def\labelenumi{\alph{enumi})}
\setcounter{enumi}{1}
\tightlist
\item
  设 n \textgreater{} 1 为一个整数，且 \(\alpha_{i}\) ( \(1 \leqslant i \leqslant n\) ) 为 C 中以 \(\delta\) 为一个最大公因子的 n 个元素。证明存在 C 上一个 n 阶方阵 A，其第一列为 \((\alpha_{i})\) 且其行列式为 \(\delta\) （利用 a））。
\end{enumerate}

\begin{enumerate}
\def\labelenumi{\arabic{enumi})}
\setcounter{enumi}{20}
\tightlist
\item
  设 \(C\) 是一个主理想整环，并设 \(A = (\alpha_{ij})\) 为一个 \(n\) 阶方阵，定义在 \(C\) 上。
\end{enumerate}

\begin{enumerate}
\def\labelenumi{\alph{enumi})}
\item
  证明存在 C 上一个 n 阶方阵 U，其行列式为 1，使得矩阵 UA 的最后一列的前 \((n-1)\) 个元素为零，且最后一个元素等于 A 的第 n 列诸元素的一个最大公因子（利用习题 20，b））。
\item
  由 a) 推出，存在 C 上一个 n 阶方阵 V，其行列式为 1，使得矩阵 \(VA = (\beta_{ij})\) 是下三角的（«A 的埃尔米特简化型»）。
\item
  若 W 是 C 上一个 n 阶方阵，其行列式为 1，且 \(WA = (\gamma_{ij})\) 是下三角的，证明 \(\beta_{ii}\) 与 \(\gamma_{ii}\) 相伴，对 \(1 \leqslant i \leqslant n\) 。
\end{enumerate}

\begin{enumerate}
\def\labelenumi{\arabic{enumi})}
\setcounter{enumi}{21}
\tightlist
\item
  一个主理想整环上的模M被称为具有有限秩，如果存在一个整数 \(m \geqslant 0\) 具有如下性质：M 的每个有限生成子模都包含在一个由至多 m 个元素生成的子模中。具有该性质的最小整数 m 称为 M 的秩。
\end{enumerate}

\begin{enumerate}
\def\labelenumi{\alph{enumi})}
\item
  证明若 \(\mathbf{M}\) 具有有限秩，则每个商模都具有有限秩，且至多等于 \(\mathbf{M}\) 的秩。
\item
  对于无挠模 M，上述定义的秩概念与 II，第 315 页中定义的概念一致。
\item
  证明若 \(\mathbf{M}\) 具有有限秩，则 \(\mathbf{M}\) 的每个子模都具有有限秩，且至多等于 \(\mathbf{M}\) 的秩。
\item
  设 \(\mathbf{M}\) 是一个 \(\pi\) -准素模，满足 \(\mathbf{M}(\pi) = \mathbf{M}\) （采用 VII，第 7 页的记号）。证明若 \(\mathbf{M}\) 具有有限秩，则它是有限生成的。
\item
  设 M 为一个没有无限高元素的 \(\pi\) -准素模。证明若 M 具有有限秩，则它是有限生成的，并且秩的概念与
\end{enumerate}

习题 8 中所定义的一致。（注意若 \(\mathbf{M}(\pi)\) 是有限生成的，则 \(\mathbf{M}\) 也是有限生成的，这可通过证明 \(\mathbf{M}(\pi)\) 的元素在 \(\mathbf{M}\) 中的高度是有界的来得到；利用 VII 第 55 页习题 5 证明这一点。）\(f)\) 设 \(\mathbf{M}\) 是一个 \(\pi\) -准素模，具有有限秩 \(r\) ，并设 \(\mathbf{M}_0\) 表示 \(\mathbf{M}\) 中由所有无限高元素组成的子模。证明 \(\mathbf{M}_0\) 是一个可除模，并且是有限个（\(p \leqslant r\) 个）不可分解子模 \(\mathrm{U}_{\pi}\) （VII，第 54 页，习题 3）的直和，且 \(\mathbf{M}\) 是 \(\mathbf{M}_0\) 与一个有限生成子模 \(\mathbf{N}\) （秩为 \(r - p\) ）的直和。（注意 \(\mathbf{M} / \mathbf{M}_0\) 是有限生成的，利用 \(e\) )；然后应用 VII 第 54 页习题 3。）\(g)\) 由上述内容推出，一个 A-模 \(\mathbf{M}\) 具有有限秩 \(r\) ，当且仅当商模 \(\mathbf{M} / \mathbf{T}\) （\(\mathbf{M}\) 对其挠子模 \(\mathbf{T}\) 的商）具有有限秩 \(p \leqslant r\) ，且每个 \(\pi\) -准素分量（属于 \(\mathbf{T}\) ）的秩都 \(\leqslant r - p\) ，并且其中至少有一个的秩为 \(r - p\) 。

\begin{enumerate}
\def\labelenumi{\alph{enumi})}
\setcounter{enumi}{7}
\item
  设 \((p_{n})\) 是互异素数的序列；设 E 为循环 Z-模 \(Z/p_{n}^{2}Z\) 的直和，并设 \(e_{n}\) 表示这样的元素：其坐标除第 n 个外均为零，而第 n 个坐标等于 \(1 \mod p_{n}^{2}\) 的剩余类。设 M 为 Z-模 \(E \times Q\) 的由元素 \((e_{n}, 1/p_{n})\) 生成的子模。证明 M 是秩为 2 的 Z-模，M 的挠子模 T 由元素 \((p_{n}e_{n}, 0)\) 生成，且 T 不是 M 的直因子。（若 \(\bar{x}_{n}\) 是 \((e_{n}, 1/p_{n}) \mod T\) 的类，证明对所有 \(m \neq n\) 存在类 \(\bar{y}_{nm} \mod T\) 使得 \(\bar{x}_{n} = p_{m}\bar{y}_{nm}\) ，但不存在元素 \(x_{n}\) 属于类 \(\bar{x}_{n}\) ，使得对所有 \(m \neq n\) 存在 \(y_{nm} \in M\) 满足 \(x_{n} = p_{m}y_{nm}\) 。）
\item
  若 \(\mathbf{M}\) 是有限生成的 A-模，证明 \(\mathbf{M}\) 的秩等于 \(\gamma (\mathbf{M})\) ，即 \(\mathbf{M}\) 的生成集的基数中最小者。
\item
  设 M 是一个有限秩的挠 A-模。证明：若 N 是 M 的真子模，则 N 不能与 M 同构（利用 f）。
\end{enumerate}

\begin{enumerate}
\def\labelenumi{\arabic{enumi})}
\setcounter{enumi}{22}
\item
  设 A 为主理想整环，且 M 为有限生成的 A-模。证明：若 N 和 P 是两个 A-模，使得 \(M \oplus N\) 与 \(M \oplus P\) 同构，则 N 与 P 同构。（化归到 M 为循环模的情形；把 \(M \oplus N\) 与 \(M \oplus P\) 等同起来，我们可以考虑一个以两种方式分解为子模直和 \(F \oplus N\) 与 \(G \oplus P\) 的 A-模 E，其中 F 和 G 是循环的且同构的。在 F 和 G 同构于 A 的情形，证明若 \(N \neq P\) ，则 N 是 \(N \cap P\) 与一个同构于 A 的子模的直和。若 F 和 G 同构于 \(A/A\pi^{n}\) ，其中 \(\pi\) 是 A 的一个不可约元，考虑 F、G 的生成元 a、b，并考察 \(\pi^{n-1}a \notin G\) 、\(\pi^{n-1}b \notin F\) 的情形，最后是 \(\pi^{n-1}a \in G\) 与 \(\pi^{n-1}b \in F\) 同时成立的情形；在最后这种情形，考虑 E 的元素 \(a + b\) 及其生成的子模。）
\item
  阿贝尔群 G 同构于某个交换域 K 的非零元素构成的乘法群 \(K^{*}\) 的一个子群，当且仅当 G 的挠子群 T 的秩 \(\leqslant 1\) （习题 22）。（关于此条件的必要性，参见 V，第 79 页，命题 2；利用同一参考文献证明该条件对 T 同构于某个域 E 的 \(E^{*}\) 的一个子群是充分的。对所有 \(\alpha \in G/T\) ，令 \(g_{\alpha}\) 为 \(\alpha\) 在 G 中的一个陪集代表元，那么，把 T 与 \(E^{*}\) 的一个子群等同起来，我们可写成 \(g_{\alpha}g_{\beta} = e_{\alpha\beta}g_{(\alpha\beta)}\) ，对所有 \(\alpha\) , \(\beta \in G/T\) ，其中 \(e_{\alpha\beta} \in E\) 。证明存在一个以 \((b_{\alpha})_{\alpha \in G/T}\) 为基的 E-代数 B，使得 \(b_{\alpha}b_{\beta} = e_{\alpha\beta}b_{(\alpha\beta)}\) 对所有 \(\alpha\) , \(\beta\) 成立，并利用 VI，第 33 页，习题 20，c) 证明 B 是一个整环。则 B 的分式域 K 满足所需的条件。）
\item
  将 VII，第 18 页的定理 1 推广到环 A 上的左模，其中 A 不一定是交换的，但 A 的每个左理想和每个右理想都是主理想（这样的环的一个例子是张量积 \(\mathbf{K}[\mathbf{X}] \otimes_{\mathbf{K}} \mathbf{D}\) ，其中 \(\mathbf{K}\) 是一个交换域，且 \(\mathbf{D}\) 是一个 K-代数，它是一个域但不一定是交换的）。
\end{enumerate}

